\section{总结与延伸}

Meta-RL 的核心价值在于实现\textbf{快速适应}：

\begin{enumerate}
    \item 多任务 RL 学习一个覆盖所有任务的策略，meta-RL 学习快速适应新任务的能力
    \item Context-based 方法通过编码历史经验实现隐式推断
    \item Gradient-based 方法（MAML）通过学习好的初始化实现快速微调
    \item 理论上 meta-RL 等价于在任务不确定的 POMDP 中做最优决策
    \item LLM 的 in-context learning 是 meta-learning 的大规模实现
\end{enumerate}

\subsection{本讲方法地图}

\begin{table}[H]
\centering
\begin{tabular}{p{0.18\textwidth} p{0.22\textwidth} p{0.22\textwidth} p{0.22\textwidth}}
\toprule
路线 & 适应机制 & 优势 & 难点 \\
\midrule
Multi-task RL & 显式 task descriptor $z$ & 结构简单，支持参数共享 & 遇到新任务时缺乏快速适应机制 \\
Black-box Meta-RL & 经验上下文 + 记忆网络 & 表达能力强，端到端统一 & 训练难、探索难、样本效率受 outer RL 约束 \\
Gradient-based Meta-RL & 学可快速微调的初始化 & 适应机制清晰 & 二阶优化开销大，RL 场景稳定性差 \\
Structured exploration (PEARL 等) & 后验采样/内在奖励/任务建模 & 更容易把探索结构显式化 & 可能引入较强假设，在部分环境中次优 \\
\bottomrule
\end{tabular}
\caption{Lecture 13 的核心方法关系}
\end{table}

\subsection{实现 Black-Box Meta-RL 的检查清单}

真正动手实现这类算法时，最容易出错的不是公式，而是 ``任务边界'' 和 ``记忆边界''。如果训练代码把 episode、task、hidden state 和 replay buffer 混在一起，最后学到的往往只是一个普通 recurrent policy，而不是能跨任务适应的 meta-policy。

\begin{table}[H]
\centering
\begin{tabular}{p{0.2\textwidth} p{0.28\textwidth} p{0.28\textwidth}}
\toprule
模块 & 常见错误 & 更稳妥的做法 \\
\midrule
Task sampling & 同一批次里的任务边界不清晰 & 显式记录 task id / task episode 边界 \\
Hidden state & 在不同任务之间错误复用隐状态 & 只在同一任务的多轮 episode 中保留记忆 \\
Reward input & 没把 reward/history 作为适应线索输入网络 & 明确把 $(s,a,r,s')$ 序列作为上下文 \\
Meta-test protocol & 测试时直接用训练任务或泄漏任务标签 & 仅给少量新任务交互数据，严格按分布外 held-out task 评估 \\
\bottomrule
\end{tabular}
\caption{Black-box Meta-RL 的实现检查清单}
\end{table}

\begin{figure}[H]
\centering
\includegraphics[width=0.9\textwidth]{slides-images/slide-39.jpg}
\caption{课程收束：今天讲 meta-RL 基础，下一讲继续 exploration 与 meta-exploration}
\end{figure}

\newpage
\subsection{为什么 exploration 被单独拿出来讲}

这节课其实已经给出答案：在普通 RL 里，探索是为了找到高奖励行为；在 Meta-RL 里，探索还要承担\textbf{识别任务}的职责。于是 ``看起来暂时没收益'' 的动作，可能恰恰是最关键的适应步骤。也正因为这种双重角色，很多 black-box meta-RL 方法虽然概念统一，但在复杂探索环境中会显著退化。

把 exploration 单独拉出来讨论的价值在于，它提醒我们不要把 Meta-RL 简化成 ``带记忆的策略网络''。真正困难的部分是：\textbf{如何让 agent 先花少量代价搞清楚自己面对的是什么任务，再把剩余预算花在正确的执行策略上。} 这也是后续从 posterior sampling、intrinsic reward 到 task dynamics prediction 一整条研究线存在的原因。

\begin{knowledgebox}{拓展阅读}
\begin{itemize}
    \item Finn et al., ``Model-Agnostic Meta-Learning for Fast Adaptation of Deep Networks (MAML),'' ICML 2017
    \item Duan et al., ``RL$^2$: Fast Reinforcement Learning via Slow Reinforcement Learning,'' ICLR 2017
    \item Rakelly et al., ``Efficient Off-Policy Meta-Reinforcement Learning via Probabilistic Context Variables (PEARL),'' ICML 2019
    \item Zintgraf et al., ``VariBAD: A Very Good Method for Bayes-Adaptive Deep RL,'' ICLR 2020
\end{itemize}
\end{knowledgebox}

\end{document}
